% !TeX root = ../组合数学笔记.tex
% 由主文档通过 \input 载入；请编译主文档。
\setcounter{chapter}{0} % 保留原书章号。
\chapter{预备知识与符号索引}
\label{app:prerequisites}

\section{逻辑、集合与有限求和}

命题 $P\Rightarrow Q$ 只说明 $P$ 足以推出 $Q$，并不自动给出逆命题。等价 $P\Leftrightarrow Q$ 需要两个方向。否定“对每个 $x$ 都存在 $y$ 使 $P(x,y)$ 成立”，应得到“存在 $x$ 使对每个 $y$ 都不成立”；量词顺序决定 $y$ 是否允许依赖 $x$。证明计数公式时，参数范围属于结论的一部分，验证有限个取值不能替代全称命题。

有限不交并的基数可相加；一般并集先变成不交分类或使用容斥。有限和换序只是在重排同一有限项集，例如 $\sum_{k=0}^n\sum_{j=0}^k a_{j,k}=\sum_{j=0}^n\sum_{k=j}^n a_{j,k}$。形式幂级数中的无限表达式则逐系数解释；当每个系数只涉及有限项时，才可以使用本书相应的换序论证。解析无穷级数需要另外的收敛条件。

归纳证明须给出起点、归纳假设和由较小参数向较大参数的推导。强归纳可一次使用所有较小情形；良序原理保证非空非负整数集合有最小元，因此最小反例法与归纳法相容。对递归算法，还须给出严格减小的非负整数参数以证明终止。

\section{行列式与 Cauchy--Binet 公式}

本书线性代数部分使用矩阵乘法、行列式的多线性与交替性，以及秩不超过行列数等基本性质。特别地，线性相关的列使行列式为零，沿仅有一个非零元的行展开可降低阶数。递推求解中使用的 Vandermonde 矩阵为 $V_{ij}=\lambda_j^{i-1}$（$1\le i,j\le r$），其行列式为 $\prod_{i<j}(\lambda_j-\lambda_i)$：行列式在两参数相同时为零，所以被每个差因子整除；双方总次数相同，比较单项式 $\lambda_2\lambda_3^2\cdots\lambda_r^{r-1}$ 的系数确定常数为一。因此不同参数保证可逆。以下公式补足矩阵树定理所需的另一项工具。

\begin{theorem}[Cauchy--Binet 公式]
若 $U$ 为 $r\times m$ 矩阵、$V$ 为 $m\times r$ 矩阵且 $1\le r\le m$，则
\[
 \det(UV)=\sum_{\substack{S\subseteq[m]\\|S|=r}}
             \det(U_{:,S})\det(V_{S,:}),
\]
其中每个子集 $S$ 的指标均按递增顺序选取。
\end{theorem}
\begin{proof}
$UV$ 的第 $j$ 列是 $\sum_{i=1}^m U_{:,i}V_{ij}$。按列多线性展开行列式，得到对 $i_1,\ldots,i_r$ 的求和，项为 $\det(U_{:,i_1},\ldots,U_{:,i_r})\prod_{j=1}^r V_{i_jj}$。若指标重复，行列式为零。对不重复指标按其集合 $S=\{s_1<\cdots<s_r\}$ 归组，写 $i_j=s_{\pi(j)}$，行列式变为 $\operatorname{sgn}(\pi)\det(U_{:,S})$。对排列 $\pi$ 求和的另一因子正是 $\det(V_{S,:})$，于是得到公式。
\end{proof}

\section{概率、符号与模型对照}

对有限样本空间，条件概率在 $\mathbb P(B)>0$ 时定义为 $\mathbb P(A\mid B)=\mathbb P(A\cap B)/\mathbb P(B)$。给定已选择的部分信息后，条件期望就是在剩余可能结果上按更新后的概率加权平均；把下一步结果分类再求平均，可恢复当前条件期望。这是第\ref{ch:probabilistic}章逐步选择方法的依据。

常用计数序列中，$D_n$ 计数错排，$S(n,k)$ 计数 $k$ 块集合分拆，$c(n,k)$ 计数 $k$ 循环置换，$B_n$ 计数全部集合分拆，$C_n$ 计数半长 $n$ 的 Dyck 路径，$p(n)$ 计数整数分拆。其空对象初值为 $D_0=B_0=C_0=p(0)=S(0,0)=c(0,0)=1$；其他边界由具体定义决定。

整数合成保留部件顺序，整数分拆只保留部件重数，集合分拆还保留块内元素标签。标号树保留每个顶点名称，无标号树按图同构取商，平面树另保留孩子次序；它们的计数不应通过未经证明的统一除法互换。普通生成函数的系数直接给对象数，指数生成函数的系数须乘相应阶乘；两者应根据组合运算而非公式外观选用。

